\documentclass{article}
\usepackage{/globals/de}
\usepackage{/globals/anki}
\ankiprefix{private-endgeräte}
\ankitopic{Private Geräte}
\ankishow
\begin{document}
\section{Umgang mit privaten Endgeräten in Firmen} 
Damit nicht jeder immer zwei Handys zur Arbeit mitnehmen muss, ein privates Handy und ein Arbeitshandy, wurden Konzepte entwickelt die es den ArbeiterInnen erlauben ein Handy zu nutzen, sowohl privat als auch für die Arbeit.

Dies ist ressourcensparend auf der Seite des Arbeitgebers, der kein, oder weniger, Geld für die Geräte der ArbeiterInnen ausgeben muss.
\ankinote{options}{Optionen}{BYOD, COPE}

\subsection{BYOD}
Beim \emph{Bring Your Own Device} Konzept ist es die Aufgabe der ArbeitnehmerIn sich ein passendes Endgerät zu besorgen und mitzubringen, wenn die Person noch keines besitzt.

Dies ist ressourcensparend auf der Seite des Arbeitgebers, der so nicht die Endgeräte aller ArbeiterInnen bezahlen muss.
\ankinote{byod}{BYOD}{Bring Your Own Device, Privathandy an Arbeitsplatz, billiger für Arbeitgeber}

\subsection{COPE}
Beim \emph{Corporate-Owned, Personally Enabled}-Konzept gehört das Gerät der Firma, die ArbeiterIn hat aber auch das Recht es privat zu nutzen.

Dies ist ressourcensparend auf der Seite der ArbeiterIn, falls diese noch kein passendes Gerät hat.
\ankinote{cope}{COPE}{Corporate-Owned, Personally Enabled, Arbeitshandy zuhause, billiger für Arbeiter}

\subsection{Probleme}
Wird ein Handy dual genutzt kann immer zu Problemen kommen, mit dem Schutz oder der Sicherheit der Daten und Informationen die mit nach hause genommen werden, Probleme mit Lizenzen (Lizenz für Büronutzung \vs Lizenz für Privatnutzung), und mentale Probleme mit der Trennung zwischen dem privaten und der Arbeit.
\ankinote{probleme}{Probleme}{Datenschutz, Lizenzen, Vermischung}
\end{document}